\pagestyle{empty}
{\BookHeadingFont\fontsize{15}{19}\selectfont\bfseries\color{BookPrimary}Διδακτικά κουτιά — v0.26}\par
{\small Επιλογή Β · αχνότερα φόντα · ενισχυμένη λάμπα · κοινό στυλ BookCallout}\par\medskip
\begin{ThinkBox}
Ξεκινώ από τον ορισμό της μετατόπισης και εξετάζω το πρόσημο του αποτελέσματος. Το θετικό πρόσημο δηλώνει κατεύθυνση προς τη θετική φορά του άξονα.
\end{ThinkBox}
\begin{AttentionBox}
Η μετατόπιση είναι διανυσματικό μέγεθος, ενώ το διάστημα είναι μονόμετρο. Δεν ταυτίζονται σε κάθε κίνηση.
\end{AttentionBox}
\medskip
{\BookHeadingFont\bfseries\color{BookPrimary}Έλεγχος διαφορετικού ύψους}\par
{\small Η θέση του εικονιδίου προσαρμόζεται στο ύψος του αντίστοιχου κουτιού.}\par\smallskip
\begin{ThinkBox}
Στην αρχή επιλέγω θετική φορά του άξονα και σημειώνω την αρχική και τελική θέση.
Στη συνέχεια χρησιμοποιώ τις δεδομένες φυσικές ποσότητες, ελέγχω τις μονάδες
και διατυπώνω την απάντηση με το κατάλληλο πρόσημο. Η ερμηνεία του τελικού
αποτελέσματος είναι εξίσου σημαντική με τον αριθμητικό υπολογισμό.
\end{ThinkBox}
\begin{AttentionBox}
Προσέχουμε το πρόσημο.
\end{AttentionBox}
\vfill
{\footnotesize\color{BookInk}Το «Συχνό λάθος» δεν έχει προστεθεί — παραμένει υπό συζήτηση.}\par
